\documentclass[11pt]{article}

\usepackage[utf8]{inputenc}
\usepackage[T1]{fontenc}
\usepackage[spanish]{babel}
\usepackage{hyperref}
\usepackage{booktabs}
\usepackage{a4wide}
\usepackage{enumitem}
\usepackage{tabularx}
\usepackage{parskip}

\pagestyle{empty}
\title{Propuesta de proyecto: Navegación de drones consciente de la creencia bajo observabilidad parcial}
\author{Mariia Osipova}
\date{Abril 2026}

\begin{document}

\maketitle
\thispagestyle{empty}

\section*{Resumen}

La pregunta en el corazón de este proyecto se cristalizó por primera vez durante una charla de Mykel Kochenderfer
en UdeSA en junio de 2025 sobre toma de decisiones secuencial bajo incertidumbre. La idea central,
que un agente debe actuar bien precisamente cuando no sabe dónde está ni qué ocurrirá después.

Mi camino hacia esta idea fue gradual. Primero exploré teoría de control:
la matemática de hacer que un sistema haga lo que uno quiere, de manera confiable y demostrable.
Como ejercicio pequeño construí un controlador PID de altitud para un cuadricóptero simulado:
el dron despega, asciende y se estabiliza a cinco metros (\texttt{[link]}).
El proyecto era simple, pero volvió tangible el desafío central:
incluso mantener una sola altitud requiere cerrar el lazo entre sensado, estimación y actuación.

Una vez que el controlador funcionó,
la siguiente pregunta natural fue: ¿cómo puedo hacer que este dron sea más autónomo?
Esa pregunta me llevó en dos direcciones a la vez. La primera fue el aprendizaje por refuerzo:
en lugar de programar el comportamiento a mano, dejar que el agente lo descubra mediante
interacción con el entorno, tomando acciones, recibiendo recompensas y aprendiendo gradualmente
una política que maximiza el rendimiento a largo plazo. La segunda fue la observabilidad parcial:
en cualquier vuelo real, el dron no conoce su estado verdadero, solo dispone de lecturas
de sensores ruidosas, retardadas e incompletas, y debe decidir qué hacer basándose en lo que cree y no en lo que sabe.
La fortaleza del RL es que puede encontrar estrategias en entornos demasiado complejos como para modelarlos a mano;
la dificultad es que, cuando los estados están ocultos, el agente debe aprender no solo qué hacer, sino también qué creer sobre dónde está.
Creo que esta tensión entre actuar y estimar está en el núcleo de lo que hace del RL bajo observabilidad parcial un problema profundo y abierto.

El disparador específico para profundizar más fue aprender sobre drones autónomos, donde un vehículo pequeño debe navegar
un entorno desconocido sin mediciones perfectas. La pregunta ``¿cómo decide?'' resultó tener una respuesta compleja.
El vehículo mantiene una distribución de probabilidad sobre los posibles estados del mundo y elige acciones que son
buenas en expectativa sobre toda esa distribución. Este marco,
el Proceso de Decisión de Markov Parcialmente Observable (POMDP),
se ubica exactamente en la intersección de lo que me resulta más atractivo:
RL como mecanismo para adquirir comportamiento, incertidumbre como restricción central y toma de
decisiones principiada como objetivo.

En una formulación estándar de RL, un Proceso de Decisión de Markov completamente observable,
el agente ve el estado verdadero en cada paso y una política sin memoria es, en principio,
suficiente. Pero un dron con sensores ruidosos nunca ve el estado verdadero.
Ve una observación que es una vista corrupta, retrasada y parcial de la realidad.
Un agente ingenuo de RL que trate esas observaciones como verdad de terreno aprenderá
una política frágil: una que funciona en condiciones limpias pero falla cuando el ruido aumenta.
La alternativa explorada en este proyecto es darle al agente acceso a una representación de su propia incertidumbre,
una creencia liviana, para que pueda aprender a actuar con cautela cuando es incierto y con
decisión cuando tiene confianza.

La pregunta central de este proyecto es:

\begin{quote}
¿Puede un agente de aprendizaje por refuerzo navegar un cuadricóptero
de manera más efectiva bajo datos de sensores ruidosos e incompletos al mantener explícitamente
una creencia sobre su propio estado, en comparación con un agente que trata las observaciones crudas
como verdad de terreno?
\end{quote}

Esta pregunta importa porque los sistemas autónomos reales rara vez operan bajo información perfecta.
Un dron que se desempeña bien solo cuando cada lectura del sensor es limpia no es lo bastante robusto
para desplegarse fuera del laboratorio. Quiero entender si, y bajo qué condiciones,
la representación explícita de la incertidumbre vuelve al aprendizaje por refuerzo más eficiente en muestras,
más estable y más transferible a entornos novedosos, y si el propio compromiso exploración-explotación
cambia cuando el agente sabe cuánto es lo que no sabe.

\newpage

\section*{Descripción general del proyecto (?)}

En términos simples, el proyecto consiste en comparar tres maneras de controlar un dron
simulado cuando los sensores no son perfectos.

Los tres enfoques son:
\begin{itemize}[leftmargin=1.8em]
    \item \textbf{PID}: un controlador clásico escrito a mano. No aprende; solo sigue reglas fijas. Noise = viento.
    \item \textbf{RL sin información extra}: un agente que aprende a volar recibiendo solo las lecturas ruidosas de los sensores.
    \item \textbf{RL con ``creencia''}: el mismo tipo de agente, pero que además recibe información sobre cuánto puede confiar en lo que está viendo.
\end{itemize}

La tarea base es sencilla: el dron tiene que volar desde un punto de origen hasta un destino fijo.
Una vez que eso funciona, se agrega ruido a los sensores progresivamente y se mide cuál de los tres enfoques resiste mejor.

La pregunta central es si ayuda que el agente sepa cuánto no sabe. Si cuando los sensores
son poco confiables el dron aprende a moverse con más cautela, eso sería evidencia
de que darle información sobre su propia incertidumbre hace una diferencia real.

%\newpage

\section*{Resultados esperados}

En esta etapa, el proyecto es una propuesta y no un estudio terminado. Esta sección describe los resultados que espero obtener al implementarlo.

\textbf{Comparación base.} Con sensores perfectos, el controlador aprendido (PPO) debería llegar a volar
al menos tan bien como el controlador manual (PID) después de suficiente entrenamiento.
Sin embargo, PID tiene una ventaja clara: no necesita entrenarse, solo ajustarse.

\textbf{Qué pasa cuando los sensores fallan.} A medida que agrego ruido a las
lecturas del sensor, espero ver cómo el rendimiento de cada enfoque se deteriora.
El controlador PID debería degradarse de forma gradual. El agente de RL sin información de
incertidumbre podría fallar de manera más brusca, especialmente si durante el entrenamiento aprendió a
asumir que los sensores son confiables.

\textbf{La idea central.} Espero que el agente que recibe información sobre su propia incertidumbre
(el ``agente con creencia'') vuele mejor cuando los sensores son ruidosos.
Si aprendió a moverse más despacio cuando no sabe bien dónde está, eso sería una señal clara de que darle
al agente información sobre cuánto confía en sus propios sensores hace una diferencia práctica.

\textbf{Generalización.} Si entrenamos al agente con distintos niveles de ruido en lugar de uno fijo,
debería adaptarse mejor a condiciones nuevas.
El agente con creencia debería generalizar aún mejor, porque tiene una forma de representar ``cuánto confío en lo que veo'' que no depende del nivel de ruido específico.

Los resultados concretos del proyecto incluirán:
\begin{itemize}[leftmargin=1.8em]
    \item curvas de aprendizaje para cada versión del agente;
    \item gráficos que muestren las trayectorias de vuelo de cada enfoque;
    \item una comparación de cómo se degrada cada uno al aumentar el ruido;
    \item una tabla que muestre qué tan bien generaliza cada agente a condiciones no vistas durante el entrenamiento;
    \item una descripción escrita de la tarea en términos formales (POMDP).
\end{itemize}

\section*{Discusión}

El objetivo de este proyecto es entender cuándo y por qué ayuda que un agente sepa cuánto no sabe.

La motivación es intuitiva: si un dron supiera exactamente dónde está en todo momento,
podría seguir reglas simples. Pero si sus sensores son ruidosos,
no puede estar seguro de su posición real, y esa incertidumbre debería cambiar cómo actúa. La pregunta es si un agente que aprende por refuerzo puede, de hecho, aprovechar esa información y si hacerlo lo vuelve más robusto.

Hay una distinción importante entre dos situaciones.
Cuando el agente tiene acceso a toda la información relevante, actuar bien es más sencillo.
Cuando solo ve una parte de la realidad, necesita rastrear lo que \emph{cree} que es verdad.
Este proyecto prueba una versión simplificada de ese rastreo y mide si hace alguna diferencia práctica.
La idea es que si el agente puede ver cuánto confía en sus propios sensores, puede aprender a ser más cuidadoso cuando esa confianza es baja.

Hay limitaciones importantes. La representación de incertidumbre que uso es una aproximación simple,
no un modelo estadístico completo. El algoritmo de RL elegido (PPO) es uno entre muchos,
y otros podrían dar resultados distintos. El simulador, aunque realista, no es idéntico a un dron real,
así que no está garantizado que lo que aprende el agente funcione igual fuera de la simulación.

Aun con esas limitaciones, el proyecto es un punto de partida concreto.
Es lo suficientemente acotado como para construirlo y analizarlo con cuidado,
y puede extenderse más adelante hacia métodos más formales, representaciones más ricas o configuraciones más complejas.
Lo veo como el comienzo de una dirección de investigación, no como un trabajo cerrado.

\newpage

\section*{Fuentes que me parecieron relevantes}

\begin{itemize}[leftmargin=1.8em]
    \item Sutton, R. S., \& Barto, A. G. (2018). \textit{Reinforcement Learning: An Introduction} (2nd ed.). MIT Press. Disponible en: \url{http://incompleteideas.net/book/the-book.html}

    \item Kochenderfer, M. J., Wheeler, T. A., \& Wray, K. H. (2022). \textit{Algorithms for Decision Making}. MIT Press. Disponible en: \url{https://algorithmsbook.com}

    \item Kaelbling, L. P., Littman, M. L., \& Cassandra, A. R. (1998). Planning and acting in partially observable stochastic domains. \textit{Artificial Intelligence, 101}(1--2), 99--134.

    \item Silver, D., \& Veness, J. (2010). Monte-Carlo planning in large POMDPs. \textit{Advances in Neural Information Processing Systems}, 23.

    \item Somani, A., Ye, N., Hsu, D., \& Lee, W. S. (2013). DESPOT: Online POMDP Planning with Regularization. \textit{NeurIPS}, 26.

    \item Panerati, J., Zheng, H., Zhou, S., Xu, J., Prorok, A., \& Schoellig, A. P. (2021). Learning to fly --- in seconds. \textit{IEEE Robotics and Automation Letters, 6}(2), 2288--2295.

    \item Schulman, J., Wolski, F., Dhariwal, P., Radford, A., \& Klimov, O. (2017). Proximal Policy Optimization Algorithms. \textit{arXiv:1707.06347}.

    \item Hwangbo, J., Sa, I., Siegwart, R., \& Hutter, M. (2017). Control of a quadrotor with reinforcement learning. \textit{IEEE Robotics and Automation Letters, 2}(4), 2096--2103.

    \item Hoel, C.-J., Wolff, K., \& Laine, L. (2020). Tactical decision-making in autonomous driving by reinforcement learning with uncertainty estimation. \textit{arXiv:2004.10439}.

    \item Tobin, J., Fong, R., Ray, A., Schneider, J., Zaremba, W., \& Abbeel, P. (2017). Domain randomization for transferring deep neural networks from simulation to real world. \textit{IROS 2017}.

    \item Kochenderfer, M. J. (2015). \textit{Decision Making Under Uncertainty: Theory and Application}. MIT Press.

    \item Mern, J., Yildiz, A., Sunberg, Z., Mukerji, T., \& Kochenderfer, M. J. (2021). Bayesian Optimized Monte Carlo Planning. \textit{AAAI 2021}.
\end{itemize}

\end{document}
